\exercisesfor{8}

\begin{enumerate}
\def\labelenumi{\arabic{enumi}.}
\item
  The mapping \(z \mapsto \bar{z}\) of \(\mathbf{C}\) into \(\mathbf{C}\) is a non-analytic continuous automorphism of the complex Lie group \(\mathbf{C}\) .
\item
  Let G be the Hilbert space of sequences \((\lambda_{1}, \lambda_{2}, \ldots)\) of real numbers such that \(\sum_{i} \lambda_{i}^{2} < +\infty\) . Consider G as a real Lie group. Let \(G_{n}\) be the set of \((\lambda_{1}, \lambda_{2}, \ldots) \in G\) such that \(\lambda_{m} \in \frac{1}{m} Z\) for \(1 \leqslant m \leqslant n\) . The \(G_{n}\) are closed Lie subgroups of G and hence \(H = \bigcap_{n} G_{n}\) is a closed subgroup of G. But this group is totally disconnected and not discrete and hence is not a Lie subgroup of G.
\end{enumerate}

¶ 3. In \(Q_{p} \times Q_{p}\) , every closed subset can be defined by a family of analytic equations. Deduce that Corollary 2 (ii) of Theorem 2 is no longer true if the hypothesis that I is finite is omitted.

¶ 4. Let G be a finite-dimensional real Lie group and A a subgroup of G. We say that an element x of L(G) is A-accessible if, for every neighbourhood U of e, there exists a continuous mapping \(\alpha\) of {[}0, 1{]} into A such that \(\alpha(0) = e\) and \(\alpha(t) \in \exp(tx)\) . U for \(0 \leqslant t \leqslant 1\) . Let b be the set of A-accessible elements of L(G).

\begin{enumerate}
\def\labelenumi{(\alph{enumi})}
\item
  Show that \(\mathfrak{h}\) is a Lie subalgebra of \(\mathbf{L}(\mathbf{G})\) . (Use Exercise 19 of §6.)
\item
  Let H be an integral subgroup of G such that \(\mathrm{L}(\mathrm{H}) = \mathfrak{h}\) . Show that \(H \subset A\) . (Let I = \(\{-1, 1\}\) and let \(R'\) be given the Euclidean norm. Let \((x_{1}, \ldots, x_{r})\) be a basis of h. Construct continuous mappings \(\alpha_{1}, \ldots, \alpha_{r}\) of I into A and continuous mappings \(f_{1}, \ldots, f_{r}\) from \(I'\) to R, such that, for \(t = (t_{1}, \ldots, t_{r}) \in I'\) ,
\end{enumerate}

\[
\begin{array}{c} \alpha_ {1} (t _ {1}) \ldots \alpha_ {r} (t _ {r}) = \exp (f _ {1} (t) x _ {1}) \ldots \exp (f _ {r} (t) x _ {r}) \\ \| t - (f _ {1} (t), \ldots , f _ {r} (t)) \| \leqslant \frac {1}{2}. \end{array}
\]

Then apply the following theorem: let \(f\) be a continuous mapping of \(\mathbf{I}^r\) into \(\mathbf{R}^r\) such that \(\| f(x) - x \| \leqslant \frac{1}{2}\) for all \(x \in \mathbf{I}^r\) ; then \(f(\mathbf{I}^r)\) contains a neighbourhood of 0 in \(\mathbf{R}^r. \dagger\)

\begin{enumerate}
\def\labelenumi{(\alph{enumi})}
\setcounter{enumi}{2}
\item
  Deduce that \(\mathfrak{h}\) is the subalgebra tangent to A at \(e\) .
\item
  Show that if A is arcwise connected, then A = H.†
\end{enumerate}

¶ 5. Let G be a Hausdorff topological group, H a closed subgroup of G and \(\pi\) the canonical mapping of G onto G/H. Suppose that H is a finite-dimensional real Lie group. There exist a neighbourhood U of \(\pi(e)\) in G/H and a continuous mapping \(\sigma\) of U into G such that \(\pi \circ \sigma = \mathrm{Id}_{\mathrm{U}}\) . (Let \(\rho\) be an analytic linear representation of H on \(\mathbf{GL}(n,\mathbf{R})\) , which is locally a homeomorphism (§ 6, Corollary to Theorem 1). Let \(f\) be a continuous function \(\geqslant 0\) on G, equal to 1 at \(e\) and zero outside a sufficiently small neighbourhood of \(e\) . Let \(ds\) be a left Haar measure of H. For \(x \in G\) , we write

\[
g (x) = \int f (x s) \rho (s) ^ {- 1} d s \in \mathbf {M} _ {n} (\mathbf {R}).
\]

Then \(g(xt) = g(x)\varphi(t)\) for \(x \in G\) and \(t \in H\) . If V is sufficiently small,

\[
g (x) \in \mathbf {G L} (n, \mathbf {R})
\]

for \(x\) sufficiently close to \(e\) . Finally use the fact that the theorem to be proved is true locally for \(\mathbf{GL}(n, \mathbf{R})\) and \(\varphi(\mathbf{H})\) .

\begin{enumerate}
\def\labelenumi{\arabic{enumi}.}
\setcounter{enumi}{5}
\item
  Let G be a group of class \(C^{r}\) (§ 5, Exercise 1) with \(r \geqslant 2\) . There exists on G one and only one real Lie group structure S such that the underlying manifold structure of class \(C^{r}\) of S is the given structure. (The uniqueness of S follows from Corollary 1 to Theorem 1. Let \(L(G)\) be the normable Lie algebra associated with G in Exercise 2 of § 5. There exist a real Lie group germ \(G^{r}\) and an isomorphism h of \(L(G^{\prime})\) onto \(L(G)\) (§ 4, Theorem 3). Verify as in § 4, no. 1 that there exist a symmetric open neighbourhood \(G^{\prime\prime}\) of \(e_{G}\) in \(G^{r}\) and a mapping \(\phi\) of class \(C^{r}\) of \(G^{\prime\prime}\) into G such that \(T_{e}(\phi) = h\) and \(\phi(g_{1}g_{2}) = \phi(g_{1})\phi(g_{2})\) for \(g_{1}, g_{2}\) in \(G^{\prime\prime}\) . By shrinking \(G^{\prime\prime}\) , it can be assumed that \(V = \phi(G^{r})\) is open in G and that \(\phi\) is an isomorphism of class \(C^{r}\) of the manifold \(G^{\prime\prime}\) onto the manifold V. There therefore exists on V a real Lie group germ structure such that the underlying manifold structure of class \(C^{r}\) is the given structure. For all \(g \in G\) , Int g defines as analytic mapping of \(V \cap (g^{-1}Vg)\) onto \((gVg^{-1}) \cap V\) (Theorem 1). By Proposition 18 of § 1, there exists on G a real Lie group structure S inducing the same analytic structure as V on an open neighbourhood of e. By translation, the underlying manifold structure of class \(C^{r}\) of S on G is the given structure.)
\item
  Let \(G\) be a real Lie group and \(H\) a closed subgroup of \(G\) .
\end{enumerate}

\begin{enumerate}
\def\labelenumi{(\alph{enumi})}
\tightlist
\item
  Let \(\mathfrak{h}\) be the set of \(x\in \dot{\mathbf{L}} (\mathbf{G})\) such that \(\exp (tx)\in \dot{\mathbf{H}}\) for all \(t\in \mathbf{R}\) . Then \(\mathfrak{h}\) is a Lie subalgebra of \(\mathrm{L}(\mathrm{G})\) . (Use Proposition 8 of §6.)
\end{enumerate}

\begin{enumerate}
\def\labelenumi{(\alph{enumi})}
\setcounter{enumi}{1}
\tightlist
\item
  Suppose that \(\mathbf{H}\) is locally compact. Show that \(\mathbf{H}\) is a Lie subgroup of \(\mathbf{G}\) . (Show first that \(\mathfrak{h}\) is finite-dimensional by proving the existence in \(\mathfrak{h}\) of a precompact neighbourhood of 0. Then imitate the proof of Theorem 2, choosing \(\mathrm{V}_2\) such that \((\exp \mathrm{V}_2) \cap \mathrm{H}\) is relatively compact.)
\end{enumerate}

